\documentclass[../main.tex]{subfiles}
\begin{document}

Chương này trình bày các công nghệ và nền tảng lý thuyết được sử dụng để xây dựng hệ
thống, kèm theo lý do lựa chọn so với các phương án thay thế. Các nội dung gồm: nền tảng
Zalo Mini App, bộ công nghệ phía giao diện, hệ sinh thái Firebase, thuật toán HMAC-SHA256,
nhận diện khuôn mặt và mã QR.

% ─────────────────────────────────────────────
\section{Nền tảng Zalo Mini App}
\label{sec:zalo}

\textit{Zalo Mini App} là nền tảng cho phép xây dựng các ứng dụng chạy trực tiếp bên
trong ứng dụng Zalo mà không cần cài đặt riêng. Về bản chất, một Mini App là một ứng
dụng web chạy trong môi trường WebView do Zalo cung cấp, được bổ sung một bộ công cụ
phát triển (ZMP SDK) để truy cập các năng lực của thiết bị và các dịch vụ của Zalo~\cite{zalominiapp}.

\subsection{Kiến trúc và mô hình hoạt động}
\label{sec:zalo-kientruc}

Về mặt kiến trúc, một Zalo Mini App là một ứng dụng web (HTML, CSS, JavaScript) được nạp và
thực thi bên trong một WebView do ứng dụng Zalo quản lý, thay vì được cài đặt như một ứng dụng
gốc. Mô hình hoạt động có thể chia thành ba lớp: (1) \textit{lớp giao diện} là ứng dụng web
viết bằng React, chịu trách nhiệm hiển thị và xử lý tương tác của người dùng; (2) \textit{lớp
cầu nối} ZMP SDK --- một lớp JavaScript trung gian chuyển các lời gọi API trong mã ứng dụng
thành lệnh tới phần gốc (native) của Zalo và trả kết quả ngược lại; (3) \textit{lớp năng lực
nền tảng} do Zalo và hệ điều hành cung cấp, gồm máy ảnh, định vị, lưu trữ cục bộ, thông tin
người dùng và các dịch vụ xã hội của Zalo. Nhờ cách tổ chức này, ứng dụng vừa tận dụng được
hệ sinh thái web (thư viện React, công cụ đóng gói) vừa truy cập được năng lực thiết bị thông
qua SDK, trong khi người dùng không phải cài đặt thêm bất kỳ phần mềm nào.

Hệ thống sử dụng các nhóm API chính sau của Zalo SDK: xác thực và lấy thông tin người
dùng (\texttt{getUserInfo}, \texttt{getPhoneNumber}, \texttt{getAccessToken}); truy cập
máy ảnh để chụp ảnh khuôn mặt và quét mã QR (\texttt{createCameraContext},
\texttt{scanQRCode}); truy cập vị trí (\texttt{getLocation}); và lưu trữ cục bộ trên
thiết bị (\texttt{setItem}, \texttt{getItem}). Các API này được nạp theo nhu cầu từ gói
\texttt{zmp-sdk}, kết hợp với bộ thành phần giao diện chuẩn của nền tảng (\texttt{zmp-ui},
hay ZaUI) để bảo đảm giao diện nhất quán với trải nghiệm của ứng dụng Zalo~\cite{zmpdocs}.

\subsection{Quy trình xây dựng và phát hành ứng dụng}
\label{sec:zalo-quytrinh}

Việc phát triển một Mini App tuân theo quy trình do nền tảng quy định~\cite{zmpdocs}, có thể
tóm lược thành năm bước từ đăng ký ứng dụng, khởi tạo mã nguồn, lập trình, chạy thử cho đến
phát hành lên hệ thống Zalo. Toàn bộ quy trình được minh hoạ trong Hình~\ref{fig:qt-taoapp}.

\begin{enumerate}[leftmargin=1.6em]
  \item \textbf{Đăng ký và tạo Mini App trên cổng nền tảng.} Nhà phát triển đăng nhập cổng dành
  cho nhà phát triển của Zalo (\texttt{developers.zalo.me}) và bảng điều khiển Mini App tại
  \texttt{mini.zalo.me}, chọn chức năng \textit{tạo Mini App mới} rồi khai báo các thông tin cơ
  bản: tên ứng dụng, mô tả và lĩnh vực (danh mục) hoạt động. Sau khi tạo, hệ thống cấp một
  \textbf{mã định danh ứng dụng (App ID)}; tại đây nhà phát triển cũng khai báo trước các quyền
  mà ứng dụng sẽ xin của người dùng (thông tin cá nhân, số điện thoại, vị trí)~\cite{zalominiapp}.

  \item \textbf{Cài đặt công cụ và khởi tạo mã nguồn.} Trên máy phát triển, cài đặt Node.js và
  công cụ dòng lệnh \texttt{zmp-cli}. Lệnh \texttt{zmp create} khởi tạo một dự án mới: công cụ
  lần lượt hỏi App ID vừa được cấp, mã truy cập (access token), khung nền tảng (đồ án chọn React)
  và mẫu khởi đầu (template), sau đó sinh ra bộ khung dự án sẵn sàng để phát triển.

  \item \textbf{Tìm hiểu cấu trúc dự án.} Mã nguồn ứng dụng được đặt trong thư mục \texttt{src},
  trong đó \texttt{pages} chứa các màn hình, \texttt{components} chứa các thành phần giao diện
  dùng lại, còn \texttt{services} và \texttt{hooks} xử lý nghiệp vụ và gọi dịch vụ. Tệp
  \texttt{app-config.json} khai báo cấu hình toàn cục của Mini App như tên ứng dụng, biểu tượng,
  màu chủ đạo và danh sách trang. Đồ án dùng Vite làm công cụ đóng gói, tích hợp với nền tảng
  qua \texttt{zmp-vite-plugin}; do vậy ứng dụng được chạy bằng \texttt{zmp start} thay vì gọi
  \texttt{vite} trực tiếp.

  \item \textbf{Bắt đầu lập trình.} Mỗi màn hình được hiện thực dưới dạng một thành phần React
  trong \texttt{src/pages} và được khai báo tuyến đường (route) tương ứng; giao diện được dựng
  bằng React kết hợp bộ thành phần ZaUI (\texttt{zmp-ui}), còn các năng lực thiết bị và dịch vụ
  của Zalo được gọi qua các API của \texttt{zmp-sdk}. Lệnh \texttt{zmp start} khởi động một máy
  chủ phát triển cục bộ, cho phép xem trước ứng dụng trên trình duyệt và tự động nạp lại khi mã
  nguồn thay đổi (hot reload); trong giai đoạn này có thể bật chế độ dữ liệu giả lập để kiểm thử
  các luồng nghiệp vụ mà không phụ thuộc vào máy chủ thật.

  \item \textbf{Phát hành.} Khi ứng dụng đã sẵn sàng, lệnh \texttt{zmp login} xác thực tài khoản
  nhà phát triển, còn \texttt{zmp deploy} đóng gói và tải bản dựng lên nền tảng. Hệ thống sinh
  ra một mã QR; quét mã này bằng ứng dụng Zalo sẽ mở Mini App ngay trên thiết bị thật, phục vụ
  việc kiểm thử trong môi trường gần với thực tế. Trước khi đến tay người dùng cuối, Mini App
  phải trải qua quy trình kiểm duyệt của Zalo nhằm bảo đảm tuân thủ chính sách của nền
  tảng~\cite{zmpdocs}.
\end{enumerate}

\begin{figure}[H]
\centering
\includegraphics[width=0.6\textwidth]{figures/qt_taoapp.png}
\caption{Quy trình xây dựng và phát hành một Zalo Mini App}
\label{fig:qt-taoapp}
\end{figure}

Các lệnh cụ thể cho từng bước cài đặt, chạy thử và triển khai được liệt kê chi tiết trong
Danh sách mã nguồn~\ref{lst:dev} và Danh sách mã nguồn~\ref{lst:deploy} ở phần phụ lục.

\subsection{Lựa chọn nền tảng và mô hình quyền}
\label{sec:zalo-lua-chon}

\textbf{Lý do lựa chọn.} So với việc phát triển ứng dụng gốc (native) hay ứng dụng web
lũy tiến (PWA), Zalo Mini App có lợi thế nổi bật là tiếp cận được lượng lớn người dùng
Zalo tại Việt Nam mà không yêu cầu cài đặt, đồng thời cung cấp sẵn cơ chế xác thực danh
tính. Bảng~\ref{tab:nentang} so sánh ba phương án.

\begin{table}[H]
\centering
\small
\caption{So sánh Zalo Mini App với PWA và ứng dụng gốc}
\label{tab:nentang}
\begin{tabular}{|p{2.9cm}|p{3.4cm}|p{3.1cm}|p{3.1cm}|}
\hline
\textbf{Tiêu chí} & \textbf{Zalo Mini App} & \textbf{PWA} & \textbf{Ứng dụng gốc} \\
\hline
Cài đặt & Không cần & Không cần & Cần cài qua store \\ \hline
Tiếp cận người dùng & Rất rộng (qua Zalo) & Trung bình & Phụ thuộc lượt tải \\ \hline
Xác thực sẵn có & Có (tài khoản Zalo) & Không & Không \\ \hline
Truy cập API thiết bị & Qua ZMP SDK & Hạn chế & Đầy đủ \\ \hline
Chi phí phát hành & Thấp & Thấp & Cao \\ \hline
\end{tabular}
\end{table}

% ─────────────────────────────────────────────
\textbf{Mô hình quyền và vòng đời.} Theo chính sách của nền tảng, Mini App không được yêu cầu
cấp quyền ngay khi mở mà chỉ xin quyền khi người dùng chủ động thực hiện thao tác liên quan;
các quyền (thông tin cá nhân, số điện thoại, vị trí) được xin theo từng phạm vi riêng biệt.
Định danh người dùng có thể lấy qua \texttt{getUserID} mà không cần quyền, thuận tiện cho việc
khởi tạo phiên đăng nhập. Nền tảng cũng cung cấp các sự kiện vòng đời (tạm dừng, tiếp tục,
thay đổi trạng thái mạng) để ứng dụng phản ứng phù hợp.

\section{Bộ công nghệ phía giao diện}
\label{sec:frontend}

Phần giao diện của ứng dụng được xây dựng trên các công nghệ sau:

\begin{itemize}[leftmargin=1.5em]
  \item \textbf{React 18}~\cite{reactjs} -- thư viện xây dựng giao diện theo mô hình thành
  phần (component), hỗ trợ cơ chế render đồng thời (concurrent) giúp giao diện mượt mà.

  \item \textbf{TypeScript}~\cite{typescriptlang} -- ngôn ngữ mở rộng của JavaScript với hệ
  thống kiểu tĩnh, giúp phát hiện sớm lỗi ngay khi viết mã và tăng độ tin cậy của hệ thống.

  \item \textbf{Jotai} -- thư viện quản lý trạng thái theo mô hình ``nguyên tử'' (atomic),
  cho phép chia trạng thái ứng dụng thành các đơn vị nhỏ độc lập. So với Redux, Jotai có
  cú pháp gọn nhẹ và phù hợp với các ứng dụng vừa và nhỏ.

  \item \textbf{Tailwind CSS} -- framework CSS theo hướng tiện ích (utility-first), kết
  hợp với SCSS để xây dựng giao diện nhất quán với bộ nhận diện của trường.

  \item \textbf{Vite} -- công cụ đóng gói (bundler) tốc độ cao, được tích hợp với plugin
  riêng của Zalo Mini App để biên dịch và đóng gói ứng dụng.
\end{itemize}

% ─────────────────────────────────────────────
\section{Hệ sinh thái Firebase}
\label{sec:firebase}

Phần xử lý phía máy chủ được xây dựng trên nền tảng Firebase của Google, gồm các thành
phần:

\begin{itemize}[leftmargin=1.5em]
  \item \textbf{Firestore}~\cite{firestore} -- cơ sở dữ liệu dạng tài liệu (NoSQL) có khả
  năng đồng bộ theo thời gian thực, phù hợp với việc cập nhật trạng thái phiên điểm danh
  trực tiếp tới thiết bị của giảng viên.

  \item \textbf{Cloud Functions} -- môi trường thực thi mã phía máy chủ không cần quản
  lý hạ tầng (serverless), được dùng để xử lý các tác vụ nhạy cảm như ký và xác minh mã
  QR, phát hiện gian lận, tránh đặt logic quan trọng phía thiết bị người dùng.

  \item \textbf{Firebase Authentication} -- dịch vụ xác thực, được dùng kết hợp với cơ
  chế cấp token tùy biến (Custom Token) để bắc cầu danh tính từ Zalo sang Firebase.

  \item \textbf{Firebase Hosting và Storage} -- lần lượt phục vụ lưu trữ và phân phối
  trang quản trị web, và lưu trữ dữ liệu ảnh.
\end{itemize}

\textbf{Lý do lựa chọn.} Firebase cung cấp gói miễn phí (Spark) đủ cho phạm vi thử
nghiệm, tích hợp sẵn khả năng đồng bộ thời gian thực và tự co giãn, giúp rút ngắn đáng kể
thời gian phát triển so với việc tự xây dựng máy chủ.

% ─────────────────────────────────────────────
\textbf{Mô hình dữ liệu của Firestore.} Firestore tổ chức dữ liệu thành các \textit{collection}
chứa các \textit{document}; mỗi document gồm các trường và có thể chứa các
\textit{subcollection} lồng bên trong. Mô hình này linh hoạt, không yêu cầu lược đồ cứng như
cơ sở dữ liệu quan hệ, nhưng đòi hỏi thiết kế cẩn thận để truy vấn hiệu quả. Việc kiểm soát
truy cập được thực hiện bằng \textit{quy tắc bảo mật} (Security Rules) --- một ngôn ngữ khai
báo quy định điều kiện đọc/ghi trên từng đường dẫn dữ liệu, áp dụng ngay tại tầng cơ sở dữ
liệu. Bên cạnh đó, Cloud Functions hỗ trợ hai dạng chính: hàm gọi trực tiếp từ ứng dụng
(callable) và hàm kích hoạt theo sự kiện (trigger) khi dữ liệu thay đổi; hệ thống chủ yếu
dùng hàm callable cho các nghiệp vụ nhạy cảm và một số hàm theo lịch cho phân tích định kỳ.

\section{HMAC-SHA256 và mã xác thực thông điệp}
\label{sec:hmac}

Mã xác thực thông điệp (Message Authentication Code -- MAC) là một đoạn dữ liệu ngắn
được sinh ra từ thông điệp và một khoá bí mật, dùng để xác minh tính toàn vẹn và tính
xác thực của thông điệp. HMAC (Hash-based MAC) là một cấu trúc MAC dựa trên hàm băm mật
mã, được chuẩn hoá trong RFC 2104~\cite{rfc2104}.

\textbf{Hàm băm mật mã.} Một hàm băm mật mã $H$ ánh xạ dữ liệu có độ dài bất kỳ thành một chuỗi
bit độ dài cố định và cần thoả mãn ba tính chất an toàn: kháng tiền ảnh (khó tìm $m$ khi biết
$H(m)$), kháng tiền ảnh thứ hai (khó tìm $m' \neq m$ sao cho $H(m') = H(m)$) và kháng va chạm
(khó tìm hai thông điệp bất kỳ có cùng giá trị băm). Hàm băm được sử dụng trong đồ án là
SHA-256 thuộc họ SHA-2, cho đầu ra 256 bit và được chuẩn hoá trong FIPS 180-4~\cite{fips180}.

Các hàm băm họ SHA-2 được xây dựng theo lối băm lặp khối (iterated construction). Một hệ quả
của lối xây dựng này là điểm yếu \textit{tấn công mở rộng độ dài} (length-extension): nếu chỉ
tính $H(K \,\|\, m)$ một cách trực tiếp, kẻ tấn công biết giá trị băm và độ dài khoá có thể suy
ra giá trị băm cho một thông điệp mở rộng mà không cần biết khoá. HMAC khắc phục điều này bằng
cách lồng hàm băm hai lần với hai lớp đệm khoá khác nhau.

Với hàm băm $H$, khoá bí mật $K$ và thông điệp $m$, giá trị HMAC được định nghĩa như sau:

\begin{equation}
\text{HMAC}(K, m) = H\big( (K' \oplus opad) \,\|\, H((K' \oplus ipad) \,\|\, m) \big)
\label{eq:hmac}
\end{equation}

trong đó $K'$ là khoá $K$ được chuẩn hoá về độ dài khối của hàm băm; $opad$ và $ipad$ là
hai hằng số đệm; $\oplus$ là phép XOR theo bit; và $\|$ là phép nối chuỗi. Trong đồ án,
hàm băm được sử dụng là \textbf{SHA-256}~\cite{fips180} do tính kháng va chạm cao và được
chuẩn hoá rộng rãi.

\textbf{Ứng dụng trong mã QR động.} Hệ thống tạo nội dung mã QR gồm các trường định danh
phiên, mốc thời gian và một giá trị ngẫu nhiên, sau đó ký toàn bộ bằng HMAC-SHA256 với
khoá bí mật chỉ được biết phía máy chủ. Nhờ tính chất của HMAC, một bên không nắm giữ
khoá bí mật không thể tạo ra mã QR hợp lệ; kết hợp với việc mốc thời gian thay đổi mỗi
30 giây, cơ chế này chống lại việc chụp và phát lại mã QR. Chi tiết thuật toán được trình
bày trong Chương~\ref{chapter:Experiment}.

\textbf{Tính an toàn.} Cấu trúc lồng hai lần hàm băm của HMAC giúp nó kháng được tấn công mở
rộng độ dài (length-extension) vốn ảnh hưởng tới cách dùng hàm băm một cách trực tiếp. Với
SHA-256 --- hàm băm cho đầu ra 256 bit, có độ kháng va chạm và kháng tiền ảnh cao --- việc dò
tìm khoá bí mật bằng vét cạn là bất khả thi trên thực tế. Do đó, kể cả khi quan sát được nhiều
mã QR hợp lệ, kẻ tấn công vẫn không thể suy ra khoá bí mật để giả mạo một mã mới. Về mặt lý
thuyết, HMAC đã được chứng minh là một hàm giả ngẫu nhiên (pseudorandom function -- PRF) và do
đó là một mã xác thực thông điệp an toàn, với điều kiện hàm nén bên dưới là một
PRF~\cite{bellare1996hmac}.

\textbf{Chống phát lại.} Mỗi mã QR mang theo một mốc thời gian và một giá trị ngẫu nhiên dùng
một lần (nonce). Phía máy chủ chỉ chấp nhận mã trong cửa sổ thời gian ngắn (90 giây) và từ
chối các nonce đã xuất hiện, nhờ đó một mã bị chụp lại sẽ nhanh chóng hết hiệu lực và không
thể tái sử dụng cho lần điểm danh khác.

% ─────────────────────────────────────────────
\section{Nhận diện khuôn mặt và FaceNet}
\label{sec:face-recognition}

Bài toán xác minh khuôn mặt (face verification) là xác định hai ảnh khuôn mặt có thuộc
về cùng một người hay không, khác với bài toán nhận dạng (identification) là tìm danh
tính trong một tập lớn. Hệ thống điểm danh sử dụng bài toán xác minh: so khớp khuôn mặt
thu được tại thời điểm điểm danh với khuôn mặt đã đăng ký của sinh viên.

Phương pháp hiện đại dựa trên mạng nơ-ron tích chập (CNN)~\cite{lecun1998gradient, he2016resnet}
để trích xuất một véc-tơ đặc trưng (embedding) có số chiều cố định từ ảnh khuôn mặt. Tiếp nối
các mô hình học sâu cho khuôn mặt như DeepFace~\cite{taigman2014deepface}, mô hình tiêu biểu là
FaceNet sử dụng hàm mất mát bộ ba (triplet loss) để học không gian đặc trưng sao cho khoảng
cách giữa các ảnh cùng người là nhỏ và giữa các ảnh khác người là lớn~\cite{schroff2015facenet}.
Việc so khớp khi đó quy về so sánh khoảng cách giữa hai véc-tơ đặc trưng với một ngưỡng
cho trước.

\textbf{Cấu trúc trích đặc trưng.} Mạng nơ-ron tích chập trích đặc trưng theo tầng: các tầng
tích chập (convolution) học các bộ lọc phát hiện đặc trưng cục bộ như cạnh, góc và kết cấu; các
tầng gộp (pooling) giảm kích thước không gian và tăng tính bất biến với dịch chuyển nhỏ; còn
các tầng phía sau tổng hợp lại thành véc-tơ đặc trưng. Các kiến trúc rất sâu với kết nối phần
dư (residual connection) cho phép huấn luyện ổn định và nâng cao chất lượng đặc trưng học
được~\cite{he2016resnet}.

\textbf{Hàm mất mát bộ ba.} FaceNet huấn luyện mạng sao cho các ảnh của cùng một người nằm gần
nhau còn ảnh của những người khác nhau nằm xa nhau trong không gian đặc trưng. Với mỗi bộ ba
gồm ảnh neo $x^a$, ảnh cùng người $x^p$ và ảnh khác người $x^n$, hàm mất mát cần tối thiểu hoá
được biểu diễn trong Phương trình~\ref{eq:triplet}:
\begin{equation}
\mathcal{L} = \sum_{i=1}^{N} \Big[\, \lVert f(x_i^a) - f(x_i^p) \rVert_2^2
- \lVert f(x_i^a) - f(x_i^n) \rVert_2^2 + \alpha \,\Big]_{+}
\label{eq:triplet}
\end{equation}
trong đó $f(\cdot)$ là véc-tơ đặc trưng (được chuẩn hoá về chuẩn đơn vị), $\alpha$ là biên
(margin) tối thiểu giữa cặp cùng người và cặp khác người, và $[z]_{+} = \max(z, 0)$. Việc tối
ưu hàm mất mát này tạo ra không gian đặc trưng mà ở đó khoảng cách Euclid phản ánh mức độ giống
nhau của hai khuôn mặt~\cite{schroff2015facenet}.

Trong đồ án, quá trình xác minh khuôn mặt được thực hiện phía thiết bị người dùng nhằm
bảo vệ quyền riêng tư của dữ liệu sinh trắc học. Hạn chế của cách tiếp cận này là chưa có
cơ chế phát hiện sự sống (liveness detection) ở mức nâng cao để chống giả mạo bằng ảnh
tĩnh; đây là một đánh đổi đã được ghi nhận và trình bày trong phần hạn chế ở
Chương~\ref{chapter:conclusion}.

\textbf{Quy trình xử lý.} Việc xác minh khuôn mặt trải qua bốn bước: (1) phát hiện khuôn mặt
trong khung hình; (2) căn chỉnh theo các điểm mốc (landmark) trên khuôn mặt; (3) trích véc-tơ
đặc trưng 128 chiều bằng mạng nơ-ron; và (4) so sánh khoảng cách Euclid giữa véc-tơ thu được
với véc-tơ đã đăng ký. Hai khuôn mặt được xem là của cùng một người nếu khoảng cách nhỏ hơn
một ngưỡng cho trước (trong đồ án là 0,6).

Đồ án sử dụng thư viện nhận diện khuôn mặt face-api.js~\cite{faceapijs} chạy trực tiếp trên
trình duyệt, với mô hình phát hiện gọn nhẹ để phù hợp với tài nguyên của thiết bị di động. Vì toàn bộ quá trình diễn
ra trên thiết bị, hệ thống chỉ lưu véc-tơ đặc trưng --- vốn không thể phục hồi ngược lại thành
ảnh --- thay cho ảnh khuôn mặt gốc, qua đó giảm thiểu rủi ro lộ lọt dữ liệu sinh trắc học.

% ─────────────────────────────────────────────
% ─────────────────────────────────────────────
\section{Mã QR và ứng dụng trong điểm danh}
\label{sec:qrcode}

Mã QR (Quick Response code)~\cite{iso18004} là một dạng mã vạch hai chiều có khả năng lưu trữ
lượng dữ liệu lớn hơn nhiều so với mã vạch một chiều, đồng thời cho phép quét nhanh từ nhiều
góc độ. Một đặc
điểm quan trọng của mã QR là cơ chế \textbf{sửa lỗi} (error correction) dựa trên mã
Reed--Solomon, cho phép khôi phục dữ liệu ngay cả khi một phần mã bị che khuất hoặc hư hỏng ở
một mức độ nhất định.

\textbf{Cấu trúc và khả năng sửa lỗi.} Một mã QR gồm các thành phần định vị --- ba ô tìm kiếm
(finder pattern) ở ba góc, cùng các mẫu căn chỉnh (alignment) và định thời (timing) --- bên
cạnh vùng thông tin định dạng/phiên bản và vùng dữ liệu. Chuẩn ISO/IEC 18004 định nghĩa bốn
mức sửa lỗi L, M, Q và H, lần lượt cho phép khôi phục khoảng 7\%, 15\%, 25\% và 30\% số từ mã
bị hỏng~\cite{iso18004}. Mức sửa lỗi càng cao thì khả năng quét trong điều kiện che khuất hay
bẩn càng tốt, nhưng dung lượng dữ liệu khả dụng càng giảm. Ở phiên bản lớn nhất (phiên bản 40),
một mã QR có thể chứa tối đa khoảng 2953 byte dữ liệu nhị phân. Nội dung mã QR của hệ thống ---
một chuỗi JSON đã ký HMAC --- chỉ dài khoảng vài trăm byte nên luôn nằm trong giới hạn này, kể
cả khi chọn mức sửa lỗi cao để tăng độ tin cậy khi quét.

Trong hệ thống điểm danh, mã QR được dùng làm phương tiện truyền tải gọn nhẹ cho hai mục đích:
truyền nội dung đã ký HMAC từ màn hình giảng viên tới thiết bị sinh viên, và trao đổi thông
tin xác minh giữa các sinh viên trong bước ngang hàng. Ưu điểm của mã QR là không đòi hỏi kết
nối trực tiếp giữa hai thiết bị --- chỉ cần một bên hiển thị và một bên dùng máy ảnh để quét
--- nên rất phù hợp với môi trường lớp học. Tuy nhiên, do mã QR tĩnh dễ bị chụp lại và chia
sẻ, đồ án kết hợp mã QR với chữ ký HMAC và cơ chế xoay mã theo thời gian tương tự mật khẩu dùng
một lần theo thời gian (TOTP)~\cite{rfc6238} để khắc phục nhược điểm này (mục~\ref{sec:hmac}).

% ─────────────────────────────────────────────
\subsection*{Kết chương}
Chương này đã trình bày các công nghệ nền tảng (Zalo Mini App, React, Firebase) cùng các
kiến thức lý thuyết cốt lõi (HMAC-SHA256, nhận diện khuôn mặt, mã QR) và lý do lựa
chọn chúng. Các công nghệ này là cơ sở để thiết kế và triển khai hệ thống được trình bày
trong chương tiếp theo.

\end{document}
